\chapter{模块、包与测试入门}
\label{ch:07-outline}

模块组织可复用定义，入口控制程序何时启动，测试固定输入与预期结果。三者配合，可以在调整代码结构时持续验证行为。


\begin{figure}[htbp]
  \centering
  \begin{tikzpicture}[
    book mindmap,
    level 1 concept/.append style={minimum size=22mm,text width=21mm}
  ]
    \node[text=white,minimum size=25mm,text width=19mm] {第 7 章}
      child[grow=163,level distance=38mm,concept color=DeepBlue] {
        node[text=white] {导入与入口}
      }
      child[grow=17,level distance=38mm,concept color=PythonYellow] {
        node {断言与参数化}
      }
      child[grow=-90,level distance=28mm,concept color=PythonBlue!22!white] {
        node {外部依赖隔离}
      };
  \end{tikzpicture}
  \caption*{模块导入、参数化测试与外部接口替身。}
\end{figure}


\section{模块导入机制}

模块是一个可导入的代码文件，包把相关模块组织到同一条导入路径下。
将复用函数放进模块，把启动动作收进 \code{main()}。
直接运行文件时，\code{__name__} 是 \code{"__main__"}；
导入时则是模块名，入口守卫便能控制启动动作。

先在本目录准备 \code{calc.py}，再由 \code{use_calc.py} 导入。
导入会执行模块顶层语句，因此连接服务和读命令参数都要找好位置。
两个模块互相导入时，先把共同使用的定义搬到独立模块，再梳理调用方向。


\begin{lstlisting}[title={\code{use_calc.py}}]
from calc import add

# 导入后只使用函数，由当前入口负责显示。
print(add(2, 3))
\end{lstlisting}


\textbf{运行核对}\enspace 预期输出：
\begin{lstlisting}[language={},numbers=none]
5
\end{lstlisting}

\par


\section[自动化测试基础]{自动化测试基础（pytest）}

测试像给工具量尺寸。输入、预期结果和失败位置都写在代码里，
改完实现再跑一次，就能看到原来的约定是否还成立。
\code{assert} 写在测试中描述预期，用户输入校验仍用正常的分支和异常。

用 \code{pytest.mark.parametrize} 指定多组输入与预期值。

在本目录运行 \code{python -m pytest -q}，三个参数组合应各跑一次。
PyCharm 可以从测试函数旁运行，也可以在测试窗口重跑失败项。


\begin{lstlisting}[title={\code{test_add_cases.py}}]
import pytest


def add(left, right):
    # 被测函数保持明确的输入与输出。
    return left + right


@pytest.mark.parametrize(
    "left,right,expected",
    [(2, 3, 5), (-2, 3, 1), (0, 0, 0)],
)
def test_add(left, right, expected):
    assert add(left, right) == expected
\end{lstlisting}


\section{Mock 与外部依赖隔离}

将 Mock 想成排练时的对手演员。外部接口回什么，由测试先约定好，
被测代码按正常结构处理它。这样能反复演练成功、超时和无效响应。

替换要发生在被测代码实际查找依赖的位置。示例把依赖作为参数传进去，
测试里可以核对返回值，也可以核对调用参数。
真服务的协议是否匹配，还要留给使用测试环境的集成验证。

\subsection*{动手｜把网络依赖换成可核对的替身}

测试状态查询时，只需核对自己的代码有没有拿对键、传对文件名。将请求函数当参数递进去，在测试里用 Mock 回一份可预期的数据，实际网络请求留给集成测试。

\begin{lstlisting}[title={\code{test_lookup.py}}]
from unittest.mock import Mock


def lookup_status(fetch, filename):
    # fetch 由调用方提供，结果仍按正常业务结构读取。
    return fetch(filename)["status"]


def test_lookup():
    fetch = Mock(return_value={"status": "done"})
    assert lookup_status(fetch, "notes.md") == "done"
    fetch.assert_called_once_with("notes.md")


# 当作脚本运行时，可以直接核对这条测试。
if __name__ == "__main__":
    test_lookup()
    print("test_lookup passed")
\end{lstlisting}

\textbf{运行核对}\enspace 直接运行本文件应打印 \code{test_lookup passed}。把业务函数里的状态键拼错会报 KeyError。保留 \code{test_} 函数命名后，也能交给 pytest 收集。

\begin{myunderstanding}
模块入口、可复用函数与外部依赖各有调用位置。先把边界画清，测试就能在合适的位置换掉依赖。
\end{myunderstanding}

\begin{pycharmpractice}
从测试旁的运行按钮执行一条参数化测试，再从测试窗口重跑失败项，核对输入与预期值。
\end{pycharmpractice}

\begin{pitfall}
脚本命名成 json.py 或 requests.py 可能遮住要导入的库。报导入错误时，先看 traceback 指向哪个文件。
\end{pitfall}

\subsection*{把两个导入入口分开跑}
\code{calc.py} 与 \code{use_calc.py} 放在同目录。分别运行这两个文件，
前者显示入口消息，后者只显示计算结果。


\begin{lstlisting}[title={\code{calc.py}}]
def add(left, right):
    # 复用逻辑位于模块顶层，启动动作位于入口守卫。
    return left + right


if __name__ == "__main__":
    print("从 calc.py 入口运行", add(2, 3))
\end{lstlisting}

\par

\subsection*{进阶边界演练}
给导入问题留三项信息，模块来源、入口文件和工作目录。导入会使用模块缓存，修改文件后若结果像没更新，先确认进程是否重新启动，再看是否读到了另一份同名模块。
